





%%%%%%%%%%%%%%%%%%%%%%%%%%%%%%%%%%%%%% FASE APP CURTA

\section{Fase de aproximação de curta distância e operação de berthing com \gls{mr}}



%%%%%%%%%%%%%%%%%%%%%%%%%%%%%%%%%%%%%% JUSTIFICATIVA

\section{Justificativa}
Instrução: justificativa tecnica para tratar o sistema como corpo rigido e dinâmico acoplado (lagrangiana energia cinetica)

\subsection{Complexidade introduzida pelo \gls{mr} em ambiente de microgravidade}

%%%%%%%%%%%%%%%%%%%%%%%%%%%%%%%%%%%%%%% REV BIBLIOGRAFICA

\section{Revisão bibliográfica}
ideia: o que outros autores fizeram com os mesmos objetivos do meu trabalho ou objetivos aproximados

%%%%%%%%%%%%%%%%%%%%%%%%%%%%%%%%%%%%%%% PROCEDIMENTO
\section{Procedimento da simulação}
ideia: (fluxograma) descrevendo o que fiz e indicando em que seções do documento estarão as formulações utilizadas

%%%%%%%%%%%%%%%%%%%%%%%%%%%%%%%%%%%%%%% FORM MATEMATICA
\section{Formulação matemática}

%%%%%%%%%%%%%%%%%%%%%%%%%%%%%%%%%%%%%%% SIMULACAO
\section{Simulação e resultados}

%%%%%%%%%%%%%%%%%%%%%%%%%%%%%%%%%%%%%%% DISCUSSÃO
\section{Discussão}
Comparação do trabalho com os trabalhos que identifiquei na revisão bibliográfica

Ressaltar as contribuições do trabalho

Ressaltar as lições aprendidas com a simulação e os resultados

Limitações do trabalho (o que faria se tivesse mais tempo)


%%%%%%%%%%%%%%%%%%%% CONCLUSÃO
\section{Conclusão}

Atingimento dos objetivos do trabalho

Resumo das contribuições

Discussão sobre relevância, originalidade, generalidade (ressaltar condições iniciais escolhidas não restringem as aplicações do trabalho), aplicabilidade

Trabalhos futuros (mesmo com mais tempo, isso está fora do escopo da pesquisa)

